\section{S1 逐项详细测试方案}
\subsection{开始执行前的统一作业步骤}
所有 CT 项共用如下流程，单项步骤如有差异，以该项内容为准：
\begin{enumerate}
\item 登记 DUT、操作员、工装、转台证书、PCB 固件及标定参数版本；核对工单与文件指纹。
\item 记录实验室温度、RH、芯片温度和开机预热时间；按照厂家手册完成配平、锁紧释放、运动区域清空、零位及急停检查。
\item 建立\textbf{新}运行目录，命令角/速率与转台反馈角/速率分开记录；保存全部原始帧及样本时间戳。
\item 先检查姿态参考、量程与坐标方向，再执行静止采样；若运动，则按“低速单轴试动、停车确认、正式试验”的顺序。
\item 所有统计优先使用\textbf{未参加拟合/调参/选模型}的 S1 数据；先检查数据有效性，再计算工程指标，最后按照已批准的产品限值签署结论。
\end{enumerate}

\testheading{CT-00}{参考基准、坐标、装夹及数据链核查}{P0／必做前置}
\textbf{目的：} 消除“坐标轴/重力符号/量程/参考台”错误导致所有后续指标虚假通过的风险。
\par\textbf{设备与输入：} DUT PCB、固定夹具、2TD-600、上位机、电子水平/参考装夹检验工具（如有）、有效计量证书及 S0 参数文件。
\par\textbf{步骤：}
\begin{enumerate}
\item 审查说明书、有效证书和设备铭牌型号一致；核对工装加 DUT 总质量不超过 $20\,\mathrm{kg}$，检查螺钉、线缆、限位、硬件急停。
\item 记录 PCB 坐标丝印朝向；建立 $S\rightarrow P\rightarrow B$ 坐标图；在低风险姿态人工确认外轴 X 与内轴 Z 的方向，以及正角反馈变化符号。
\item 在台面朝上静止，记录外/内轴真实反馈位置、ACC 三轴符号与模长、GYRO 三轴静止输出；适当小角度正负试动确认转台模式与停机反馈。
\item 读回 ACC/GYRO 量程、ODR、单位、MCU/固件版本以及内参参数哈希；若量程源自软件默认回退而非固件证据，标为\texttt{BLOCKED}。
\item 分别断开一次通讯或关闭记录\emph{（转台完全停止后）}，确认采集程序明确记录失联/丢帧，不沿用旧缓存。
\end{enumerate}
\textbf{输出：} \texttt{reference\_check.csv}、坐标系示意、设备证书与版本快照、正负方向核对签名。
\par\textbf{判定：} 轴系与参考定义明确、量程和补偿参数对应、设备安全项通过时，允许进入后续 CT。任何一项不明确均暂停涉及该参考量的测试。\textbf{注意：}“内轴 360° 与 0° 在反馈数值上等价”不等于两个机械/线缆路径都安全。

\testheading{CT-01}{标定参数生效与补偿前后对比}{P0／直接可做}
\textbf{目的：} 确认实际参与计算的补偿参数正是 S0 的冻结结果，并证明补偿没有被漏用、重复应用或单位错误。
\par\textbf{设备与工况：} 选择 S1 独立的至少 6 个静态姿态以及 1 个静止零速姿态；参数版本保持不变。
\par\textbf{步骤：}
\begin{enumerate}
\item 记录矩阵、零偏、坐标系、数值精度及 CRC/哈希；比较原始码值换算与工程软件导出的补偿结果。
\item 采集同一时间段的原始六轴数据与补偿六轴数据，\textbf{不要}用不同采样批次做改善率比较。
\item 在 PC 离线执行 $\hat{\bm f}_P=\mathbf C_a(\bm f_\mathrm{raw}-\bm b_a)$ 和 $\hat{\bm\omega}_P=\mathbf C_g(\bm\omega_\mathrm{raw}-\bm b_g)$；与设备或上位机实际输出逐样本核对。
\item 对独立姿态计算补偿前/后的矢量 RMSE、峰值误差、静止角误差；明确补偿前原始传感器坐标是否已换到同一参考系，否则改善率\emph{不可直接对比}。
\item 仅在安全且不会改变生产校准文件的测试镜像中做“开/关补偿”对比；不要在正式参数区反复擦写。
\end{enumerate}
\textbf{计算：} $\eta=(E_\mathrm{before}-E_\mathrm{after})/E_\mathrm{before}\times100\%$，同时报告两项绝对误差；如果分母接近 0，改善率记为“无定义”。
\par\textbf{判定：} 逐样本校正数值与冻结公式一致；改善趋势正常且独立指标满足\pending。若算法并未输出补偿六轴，保留离线验证结果，标记为“PC 离线补偿”。

\testheading{CT-02}{零位误差与静态零速残余}{P0／直接可做}
\textbf{目的：} 分别验证 PCB/平台零位下的静态倾角和零角速度；两者不可混成一个“零偏”指标。
\par\textbf{工况：} 外轴 $0^\circ$、内轴 $0^\circ$，已知夹具安装和参考零位；室温热稳态后采 3 段，每段 30\,s。
\par\textbf{步骤：}
\begin{enumerate}
\item 预热后确定实际参考零位并记录反馈角度、安装零位修正值及水平判据。
\item 取三段静止六轴数据，应用冻结 ACC/GYRO 参数；以固定测量窗分别求均值、标准差和最大值。
\item 用式\,\eqref{eq:tilt} 求静态 Roll/Pitch；分别与由转台实测位置算出的 $\theta_{\rm ref},\phi_{\rm ref}$ 比较。
\item 计算每段三轴 GYRO 平均残差 $\bar{\bm\omega}$，记录 $\|\bar{\bm\omega}\|$ 以及随时间变化。
\item 重新回到该零位再采一段，区分一次启动偏置和位置返回误差。
\end{enumerate}
\textbf{计算：} $e_{\phi,0}=\bar\phi-\phi_{\rm ref}$，$e_{\theta,0}=\bar\theta-\theta_{\rm ref}$，$b_{g,j}^\mathrm{res}=\overline{\hat\omega_j}$。
\par\textbf{判定：} 零位、零速的限值分别为\pending；不把夹具初始倾斜简单归罪为 ACC 内参错误。需要绝对水平精度时补充夹具/台基标定不确定度。

\testheading{CT-03}{ACC 三轴矢量独立验证}{P0／直接可做}
\textbf{目的：} 检验完整 3$\times$3 校正后在\textbf{未参与模型拟合或选择}的姿态上的矢量残差，而非只看 $\lVert \bm f\rVert=1$。
\par\textbf{工况：} 推荐 12--18 个安全且覆盖半球/全空间的封存姿态；要求正负与倾斜中间角兼有。与 S0 已用于选模的 12 个验证位置区分；工况开始前检验 $\mathbf H=[\bm f_{P,i}^{\mathsf T},1]$ 的秩和条件数。
\par\textbf{步骤：}
\begin{enumerate}
\item 使用已核准的姿态列表，逐点由实际转台反馈 $\alpha_i,\beta_i$ 和式\,\eqref{eq:refacc} 计算比力参考 $\bm f_{P,i}$。
\item 两轴停止、达到姿态与温度稳定条件后每点采集 10--30\,s；建议重复 3 轮并改变角度路径。
\item 先用冻结矩阵对全量数据统一补偿，再用预先登记的静态窗口提取均值；保留被剔除样本和原因。
\item 求每点三轴误差向量 $\bm e_i=\bar{\bm f}_{P,i}^{\rm corr}-\bm f_{P,i}$、各轴误差、欧氏模长、总体 RMSE 和最大值。
\item 绘制误差随真实姿态、各轴分量、测试顺序及温度的关系图；对异常方向复测，但第一次数据不得从报告中删除。
\end{enumerate}
\textbf{计算：} $\mathrm{RMSE}_{\rm vec}=\sqrt{N^{-1}\sum_i\lVert\bm e_i\rVert^2}$，$E_{\max}=\max_i\lVert\bm e_i\rVert$，单位 g。
\par\textbf{判定：} 各轴/合成误差满足\pending；数据集选择必须封存，不能因误差高而调整角度或重新拟合 C 矩阵后仍称“独立”。

\testheading{CT-04}{ACC 模长、正反向对称与非线性诊断}{P0／直接可做}
\textbf{目的：} 识别“矢量方向不对但模长接近 1g”、正负响应不对称及中间角残差放大的问题。
\par\textbf{工况：} 从 CT-03 独立姿态中挑选理论上互为 $+\bm u$ 与 $-\bm u$ 的姿态配对；如实际路径不可达则删去该配对并说明。再选若干中间角。
\par\textbf{步骤：}
\begin{enumerate}
\item 在各配对姿态下记录实际反馈角和补偿后的加速度三轴均值，而不是依据命令位置判断是否严格对向。
\item 计算模长误差 $\delta_{m,i}=\lVert\bar{\bm f}^{\rm corr}_i\rVert-1$、方向误差 $\gamma_i$。
\item 对参考方向严格相反的成对数据计算 $\bm d_{\pm}=(\bar{\bm f}^{\rm corr}_{+}+\bar{\bm f}^{\rm corr}_{-})/2$；不严格相反时，用实际参考矢量修正后再比较残差。
\item 将各误差对照倾角和输入分量分箱，分析是否在接近大倾角/某特定轴时明显增大。
\item 向 CT-03 报告增加模长、方向、对向差三项而不是只发布一条合成分数。
\end{enumerate}
\textbf{计算：} $\gamma_i=\arccos\big((\hat{\bm f}_i\cdot\bm f_{P,i})/(\|\hat{\bm f}_i\|\|\bm f_{P,i}\|)\big)$，结果转换为角度。
\par\textbf{判定：} 模长、方向及对向差限值\pending；本项只支持\emph{常温 $\pm1g$ 重力范围}下的评价，不能证明高于 $1g$ 的 ACC 全量程线性度。

\testheading{CT-05}{单方向静态倾角精度}{P0／DUT 有角度输出或 PC 离线}
\textbf{目的：} 把 ACC 矢量残差转换成使用者能理解的静态 Roll/Pitch 误差，但不与动态融合精度混淆。
\par\textbf{工况：} 按产品名义倾角范围选择\emph{可达且安全}的参考点，例如 $0^\circ,\pm1^\circ,\pm5^\circ,\pm15^\circ,\pm30^\circ,\pm45^\circ$；不应默认 PCB 具备 $\pm60^\circ$ 或更大量程。每个角度两轴反馈都必须被记录。
\par\textbf{步骤：}
\begin{enumerate}
\item 外轴改变倾斜角、内轴取预审的固定角，先确认运动学计算生成需要的 Roll/Pitch 变化和符号。
\item 按不同倾角从小到大再反向采集；每点采集三段，间隔重复到位。
\item 参考角度由\emph{真实 $\bm f_P$ 参考向量}按照与产品一致的 Roll/Pitch 定义计算；禁止把外轴命令数值直接写成 X/Y 两个倾角参考。
\item 用冻参 ACC 计算静态倾角，或保存 MCU 输出的角度；两者分别出报告。对接近 $\pm90^\circ$ 的 Euler 奇异区域谨慎使用角差，必要时仅比较矢量夹角。
\item 逐点计算两角偏差、RMSE、P95（如果样本量足够）和峰值；按正负方向单独画曲线。
\end{enumerate}
\textbf{计算：} $e_{\phi,i}=\hat\phi_i-\phi_{\rm ref,i}$、$e_{\theta,i}=\hat\theta_i-\theta_{\rm ref,i}$，对角度周期做合理归一化。
\par\textbf{判定：} 每轴角度误差与角度测量范围\pending；若只有离线结果，结论写为“\textbf{基于常温 ACC 的静态计算角度精度}”，不写“整机动态倾角精度”。

\testheading{CT-06}{复合姿态与交叉轴误差}{P0／直接可做}
\textbf{目的：} 检查在真实 3D 参考重力方向变化时，X/Y 倾角输出和轴间耦合是否正确。
\par\textbf{工况：} 审查安全后选择 $(\alpha,\beta)$ 的 8--12 个\emph{复合姿态}，如外轴 $15^\circ,30^\circ,45^\circ$ 与内轴若干非零方位角的组合；优先复用 CT-03 新采集点但额外输出角度耦合报告。
\par\textbf{步骤：}
\begin{enumerate}
\item 用 $R_x(\alpha)R_z(\beta)$ 计算每个工况的 $\bm f_{P}$，进一步求 $\phi_{\rm ref},\theta_{\rm ref}$；先审查是否覆盖目标二维角度区域。
\item 固定外轴在非零倾斜，改变内轴方位，得到\emph{坐标分量变化}的不同组合；内轴为 0 外轴为 0 时不能把内轴转动写成独立倾角激励。
\item 每个组合点到位稳定后采样 10--30\,s，保存实际位置；换向再进行一轮。
\item 分别计算 X、Y 角度残差和残差之间的相关性；检验一个参考角变化、另一个参考角应接近不变的工况是否确实成立，不能凭单轴转动自行假定。
\item 检查异常耦合是否来自软件参考模型、夹具方向、ACC 非正交残差或 Roll/Pitch 角度定义的非线性。
\end{enumerate}
\textbf{输出：} $(\phi_{\rm ref},\theta_{\rm ref})$ 空间的角误差热图、交叉轴残差曲线、异常路径说明。
\par\textbf{判定：} 复合角度残差限值\pending。与 CT-05 结果共同评价；\textbf{不以“另一轴读数为 0”作为所有工况的默认判据}。

\testheading{CT-07}{同姿态重复性}{P0／直接可做}
\textbf{目的：} 确定同一机械姿态下，测量结果在多次独立到位后的离散性。
\par\textbf{工况：} 选 $0^\circ$、中倾角、大倾角三类已核准姿态，每一姿态至少 5 次重新到位；每次采 10--30\,s。
\par\textbf{步骤：}
\begin{enumerate}
\item 每次从同一安全起点离开目标、再按一致路径返回；不要只对一段连续时间采样分成 5 份冒充重复定位。
\item 记录每次真实反馈角、温度、平均六轴与平均 Roll/Pitch；统计组内噪声与组间到位离散度。
\item 用\emph{同一个冻参}计算各次均值与总体均值的差；必要时对参考姿态的微小反馈差进行修正。
\item 对不同方向、不同倾角分别报告标准差和极差；检查是否有工装松动或 PCB 翘曲。
\end{enumerate}
\textbf{计算：} $s_{\rm rep}=\sqrt{\sum_{i=1}^{n}(\bar\theta_i-\bar{\bar\theta})^2/(n-1)}$；另报 $\max\bar\theta_i-\min\bar\theta_i$。
\par\textbf{判定：} 重复性限值\pending。\textbf{角度测量重复性}应与\textbf{转台回位重复性}分别估计或在总预算中说明不可分离部分。

\testheading{CT-08}{正反路径回程差}{P0／直接可做}
\textbf{目的：} 检查相同反馈姿态从不同方向进入时，PCB/装夹/转台及算法是否存在路径依赖误差。
\par\textbf{工况：} 同一目标角（如 $0^\circ,15^\circ,30^\circ$），分别从较小角、较大角方向接近；每目标至少 5 对试验。
\par\textbf{步骤：}
\begin{enumerate}
\item 在机械许可范围内制定往返路径，运动中单轴驱动且遵守 2TD-600 限制。
\item 每次目标到位后等待相同时间并采集相同时间长度，保存最后一段实际轴位置均值。
\item 以相同实际参考姿态计算误差：$e_{+}$ 为正向接近误差，$e_{-}$ 为反向接近误差。
\item 统计 $H=|\bar e_{+}-\bar e_{-}|$，并报告每次完整对照数据；若两方向真实位置不相同，不可直接把 $H$ 归为 DUT 的磁滞或迟滞。
\item 若发现显著回程差，重复检查装夹、驱动停止状态、轴承回程/编码器反馈和内参，而不是直接重新标定抵消。
\end{enumerate}
\textbf{判定：} 回程差指标\pending；与 CT-07 的组间重复性同时判断。

\testheading{CT-09}{GYRO 静止零速残差及姿态相关性}{P0／直接可做}
\textbf{目的：} 验证陀螺仪常温零偏校正是否有效，以及静止姿态变化造成的输出差异（可能存在重力敏感性）。
\par\textbf{工况：} 零位和至少 5 个姿态分散点；每姿态静止采 60--120\,s，最后附一个返回初始姿态工况。
\par\textbf{步骤：}
\begin{enumerate}
\item 两轴确认完全停止且状态反馈新鲜，避免将转台爬行当作 GYRO 零偏。
\item 采集三轴补偿后的静态角速度，计算每姿态均值、标准差、最大零速合矢量误差。
\item 根据姿态计算平台重力矢量，检查 $\bar{\bm\omega}$ 是否随重力方向系统变化；如变化明显，另行分析 GYRO 的 $g$ 敏感性或热/装夹干扰。
\item 重回初始姿态并重复测量，判定是否是时间漂移而非姿态相关性。
\item 报告每轴结果和 $\|\bar{\bm\omega}\|$，且明确未完成温补场景。
\end{enumerate}
\textbf{计算：} $\bm b_{g,i}^{\rm res}=\frac{1}{n_i}\sum_k\hat{\bm\omega}_{i,k}$，$B_{g,\max}=\max_i\|\bm b_{g,i}^{\rm res}\|$。
\par\textbf{判定：} 零速残余限值\pending。此项只估计\emph{静止工况零偏}，不能证明高速比例因子正确。

\testheading{CT-10}{GYRO 多速率正反向残余}{P0*／当前可做诊断，绝对精度有条件}
\textbf{目的：} 评估标定后的陀螺仪在低速正反方向下的响应、候选比例因子残余和方向不对称。\textbf{计量警告：} 2TD-600 手册没有给出角速率精度保证，因此本项绝对比例因子只有在补充有效参考速率计量依据后才能签 PASS。
\par\textbf{工况：} 采用程序当前候选速率集合：拟合集 $\pm1,\pm2,\pm5^\circ/\mathrm s$；\textbf{S1 新增验收工况}可从 $\pm1.5,\pm3.5^\circ/\mathrm s$ 中抽取，但若原 S0 已把这些速率用于模型选择，则还需\emph{另外的新速率点}及独立运行。例如试运行 $\pm0.5,\pm4^\circ/\mathrm s$，须先审核软件安全门槛与稳速能力。
\par\textbf{步骤：}
\begin{enumerate}
\item 审查安全转动路径与连续运行可用的稳态行程，先以较小速率做单轴正反运动，确认速度符号及停车逻辑。
\item 对 Z、X、Y 三个参考方向，按现行候选运动学设计可行的\emph{固定姿态+单轴匀速}方案，记录实际两轴反馈位置与速率；不能把某轴的命令速率直接当成平台坐标速率向量。
\item 在稳定匀速窗口同步采集六轴与参考反馈。通过平台角速度映射 $\bm\omega_P=\bm R_{BP}^{\mathsf T}(\dot\alpha\bm e_x+\dot\beta\bm R_x(\alpha)\bm e_z)$ 得到候选参考。
\item 分别计算每点修正后 GYRO 平均矢量残差、顺时针/逆时针均值对称性、速率和温度漂移，以及轨迹周期内是否有丢帧。
\item 若参考角速度没有计量保障，仅输出\textbf{相对线性/方向一致性诊断}；获取证书或额外标准速率后再估计比例因子误差和扩展不确定度。
\end{enumerate}
\textbf{计算：} $\bm e_{\omega,i}=\hat{\bm\omega}_{P,i}-\bm\omega_{P,i}^{\rm ref}$；参考方向单位矢量 $\bm u_i$ 上的轴向误差为 $e_{\parallel,i}=\bm u_i^{\mathsf T}\bm e_{\omega,i}$。
\par\textbf{判定：} 速率绝对误差与比例因子限值\pending。现有手册条件下仅能作\texttt{DIAGNOSTIC}；缺乏必要参考量时记为\texttt{BLOCKED}，不能判定绝对速率精度合格。

\testheading{CT-11}{GYRO 轴间耦合、正反对称与低速线性}{P0*／有条件}
\textbf{目的：} 验证六轴矩阵对非目标轴响应的抑制、速率曲线非线性和正反向响应。
\par\textbf{工况：} 使用 CT-10 的安全低速平台参考方向，并在独立稳态段复测三轴候选方向。
\par\textbf{步骤：}
\begin{enumerate}
\item 对每个角速度工况使用真实运动学参考 $\bm\omega_P^{\rm ref}$，取单位方向 $\bm u=\bm\omega_P^{\rm ref}/\|\bm\omega_P^{\rm ref}\|$。
\item 计算沿参考方向的测量值 $\omega_\parallel=\bm u^{\mathsf T}\hat{\bm\omega}_P$ 及横向向量 $\bm e_\perp=\hat{\bm\omega}_P-\omega_\parallel\bm u$。
\item 对同一参考方向上的多个速度点，拟合 $\omega_\parallel=k\omega_{\rm ref}+b$ 并报告残差；\textbf{仅在参考速率已标定}时将 $k-1$ 称为绝对比例因子误差。
\item 比较正反转结果，分别保存两方向的稳态平均值和波动幅度；检查是否有加减速段混入。
\item 记录横向响应 $\|\bm e_\perp\|$ 随速度幅值的趋势，分析多轴矩阵是否带来方向偏差。
\end{enumerate}
\textbf{计算：} $E_\mathrm{cross}=\max_i\|\bm e_{\perp,i}\|$，$\delta_k=(k-1)\times 10^6\,\mathrm{ppm}$；$\delta_k$ 须满足参考速率计量前提。
\par\textbf{判定：} 横向误差、非线性及正反差\pending。若转台轴方向未核准、校准速度无溯源，只发布算法诊断，不宣称单轴标度精度。

\testheading{CT-12}{常温静态噪声与短期漂移}{P0／直接可做}
\textbf{目的：} 在标定系数冻结后表征\emph{实际输出波动和漂移}，避免用一次短静态均值替代稳定性评估。
\par\textbf{工况：} 固定不动、外界振动尽量低；预热后连续采样建议 $30$--$60\,\mathrm{min}$，同时记录室温、芯片温度、电源。长时间 Allan 方差测量单列为 CT-12A，不在此项混用不同采样窗口。
\par\textbf{步骤：}
\begin{enumerate}
\item 保持转台完全静止和同一装夹，不在测试期间执行额外运动或更新零偏参数。
\item 保存\textbf{原始逐样本}六轴输出及实际时间戳，并分别保存经过补偿/滤波的输出；对两条链路分别评价。
\item 按 1\,s、10\,s、60\,s 等固定分组计算均值、标准差、峰峰值、角度变化范围；记录窗口长度，避免不同窗口结果混用。
\item 执行趋势分解：平均值随时间、芯片温度与室温变化的相关性；留意自热、供电波动、串口丢帧造成的假漂移。
\item 需要时使用等间隔有效采样的 Allan deviation (ADEV) 求角随机游走和零偏不稳定性；对不等间隔、缺帧和预滤波数据需明确局限。
\end{enumerate}
\textbf{计算：} 每轴 $s_x=\sqrt{\sum(x_i-\bar x)^2/(n-1)}$，$\Delta_{\rm drift}=\max_t\bar x(t)-\min_t\bar x(t)$。角度亦同。
\par\textbf{判定：} 角度波动、静态零偏和漂移限值\pending。噪声\emph{表征}不是可由固定内参矩阵完全消除的偏差。

\testheading{CT-12A}{Allan 方差／Allan 偏差与随机稳定性分析}{P1／推荐常温专项；独立连续长时记录}
\textbf{目的：} 对标定后的 ACC/GYRO 常温随机噪声进行独立表征；重点观察陀螺仪角随机游走（ARW）与零偏不稳定性（bias instability），以及可能出现的长时间尺度随机噪声项。\textbf{本项不是再次求解确定性内参矩阵，也不自动改善噪声。}
\par\textbf{设备与前置：} IMU PCB、稳定供电与连续采集程序、固定静止的已调平 2TD-600 或其他低振动固定平台；常温环境和芯片温度应同步记录。必须保持 S0 冻结参数、量程和 ODR 一致，关闭或明确记录动态自校零/自适应滤波。
\par\textbf{建议工况：} 热稳定后\textbf{连续采集 3--12 小时作为研发起点}，以实际目标的 $\tau_{\max}$ 和数据稳定性确定最终时长；优先保留原始等时间间隔的 GYRO 三轴角速度（如 $^\circ/\mathrm s$）和 ACC 三轴加速度（g），以及时间戳、数据序号、电源和温度。采样率以真实固件输出为准，不能事后虚构等间隔时间戳。
\par\textbf{详细步骤：}
\begin{enumerate}
\item 记录 DUT、固件、采样率（实测 ODR）、滤波链、量程、补偿参数哈希、安装姿态、环境温度、芯片温度、供电；确认转台静止，并记录开始和结束的参考反馈。
\item 预热到既定稳态条件后开启新的连续数据文件；全程\textbf{不复位、不转动、不重标、不切换滤波器或零偏估计模式}，每次中断需新建运行编号而不是拼接成一条连续序列。
\item 校验帧序号、时间戳单调性、丢帧率、异常值、采样间隔分布。对有缺帧或不等间隔的数据，不得直接用等间隔 ADEV 公式强行计算；可先使用完整连续段，记录长度与丢失比例。
\item 画出原始与补偿后三轴角速度随时间的曲线、芯片温度随时间的曲线，并检查室温渐变、供电扰动、轴系微动造成的非平稳趋势；不允许不说明就去趋势以获得更漂亮的 Allan 曲线。
\item 对每一轴使用\textbf{overlapping Allan deviation}（或报告明确说明非重叠算法），在合适的 $\tau$ 网格上计算 $\sigma_A(\tau)$；逐点记录有效簇数量，尾部有效对数太少的点不得用于定量结论。
\item 在 $\log\tau$--$\log\sigma_A$ 图中识别白噪声的 $-1/2$ 斜率段、近似平坦区和可能的 $+1/2$ 段；依据统一单位和明确采用的噪声模型推算 ARW、零偏不稳定性及可识别的其它系数，不应将任一单独的最小值不加说明地当成 bias-instability 规格。
\item 在报告中并列输出：原始六轴、补偿六轴（以及如有的滤波六轴）的 ADEV 图；各系数、选用的 $\tau$ 拟合范围、拟合斜率、温度趋势与评估可信度。对不同滤波链生成\textbf{分别标注}的结果，不混合评价。
\end{enumerate}
\textbf{计算：} 相邻簇均值差的 Allan 偏差定义和 $\tau$ 取值方法见附录《计算公式》；\textbf{角随机游走的单位}应报告为 $^\circ/\sqrt{\mathrm h}$（或等效标准单位），零偏不稳定性可报 $^\circ/\mathrm h$，给出换算和模型约定。
\par\textbf{输出：} \texttt{CT12A\_allan\_raw.csv}、采集质量表、每轴 ADEV 数据表与双对数曲线、三轴 ARW / bias instability 估计表、算法/温度条件、排除规则和结果可信等级。
\par\textbf{判定：} 当前先形成\texttt{CHARACTERIZATION/DIAGNOSTIC} 噪声特性基线；若未来已审批产品噪声指标、最小数据时长和统计方法，才能转为量化 PASS/FAIL。静态采集时长不足或样本非连续时，相关长时间尺度结果必须标为\texttt{BLOCKED}，不得补齐想象数据。

\testheading{CT-13}{多次上电重复性}{P0／直接可做}
\textbf{目的：} 验证同一块 PCB 参数固定时，不同启动过程中的零偏和静态角度是否一致。
\par\textbf{工况：} 夹具不移动，保持常温稳定，建议完整电源循环 $10$ 次；记录断电时长、预热时间及芯片温度。
\par\textbf{步骤：}
\begin{enumerate}
\item 对每次上电用同一个固件与冻结内参启动；在固定的预热时刻（例如 1、5、15\,min）采一段静止数据。
\item 记录每次上电 GYRO 平均零速、ACC 静态矢量、静态 Roll/Pitch、参数 CRC 和量程读回；确认没有自动零偏重估逻辑悄然变化。
\item 比较相同预热时间的跨次标准差和最大极差；若芯片自热路径不同，应分层统计并记录原因。
\item 如 DUT 设计包含开机自校零，则\textbf{分别报告}自校零前/后的结果，否则无法评价固定内参贡献。
\end{enumerate}
\textbf{计算：} $s_{\rm on}=\mathrm{std}\{\bar\theta_1,\ldots,\bar\theta_{10}\}$；GYRO 各轴同理。
\par\textbf{判定：} 跨上电最大偏差、标准差、参数加载一致性\pending；参数未加载/加载错误直接判\texttt{FAIL}（软件一致性项）。

\testheading{CT-14}{重复装夹与复测}{P1／有条件}
\textbf{目的：} 识别裸 PCB/夹具重复装夹对最终倾角的贡献，并与内参自身重复性区分。
\par\textbf{工况：} 使用可复位且受控的 PCB 工装；推荐连续 3--5 次卸装、复装。若 PCB 不适合反复拆卸则取消该项并写明原因。
\par\textbf{步骤：}
\begin{enumerate}
\item 首先完成 CT-07“不卸装”的基线重复性试验。
\item 停机断电且固定转台轴系后卸下 PCB，再按统一扭矩、定位销和夹紧流程复装；禁止在驱动运动时操作。
\item 每次复装仅采集 $0^\circ$ 和两个非零已知参考姿态，使用\textbf{同一}已冻参计算角度误差。
\item 将装夹重复性与 CT-07 同位置复测结果比较；如果复装会改变 $S\rightarrow P$ 外参，保留原始值和独立测得的外参值。
\end{enumerate}
\textbf{判定：} 复装导致的总误差\pending。此项结果体现工装与装配，不宜反向归类成 IMU 芯片内参误差。

\testheading{CT-15}{单次常温试验内的温度稳定性}{P1／直接可做}
\textbf{目的：} 确保本轮标定后精度不是由 PCB 自热和试验室温波动掩盖或放大。
\par\textbf{工况：} 室温控制在手册允许的 $20$--$30^\circ\mathrm C$ 范围，选不同预热时刻和 CT-03 的若干重复姿态；持续记录芯片温度。
\par\textbf{步骤：}
\begin{enumerate}
\item 从冷机开始连续记录 $30$--$60\,\mathrm{min}$，确认芯片温升曲线；对比开始和热稳态阶段的 GYRO 零偏及 ACC 静态角。
\item 在多个时间窗口分别计算与同一真实静态参考角的偏差；拟合仅用于\emph{诊断}的误差-芯片温度趋势。
\item 记录室温、芯片温度及二者差值，不能将芯片温度漂移误判为实验室温度超标。
\item 标注 S0 标定实际采用的温度区间；若 S1 与 S0 温度差明显，先分析是否需独立温补，而不是修改常温参数后仍称同一次 S1 验收。
\end{enumerate}
\textbf{输出：} 温度-时间图、误差-温度图、预热建议窗口。
\par\textbf{判定：} 本项以\texttt{DIAGNOSTIC} 为主；全温温漂/温补验收属于阶段 S2。

\testheading{CT-16}{低速运动后的静止恢复与回零}{P1／可做基本诊断}
\textbf{目的：} 检查标定后 GYRO/ACC 和可选姿态融合算法对低速运动停止的响应，检测运动后零点跳变。
\par\textbf{工况：} 单轴安全的往返运动，例如小幅 $\pm5^\circ$ 位置变化或短时低速正反运动（限软件试运行 $\le10^\circ/\mathrm s$）；\textbf{不做未经审核的双轴同时连续旋转}。
\par\textbf{步骤：}
\begin{enumerate}
\item 静止基线采 30\,s；记录转台实际零位与 PCB 静态角度。
\item 执行单轴低速运动—停止—返回已知参考姿态，必须确认反馈速度归零且到位。
\item 返回后在 0--5\,s、5--30\,s、30--60\,s 三个时间窗采样；比较静态角、GYRO 零偏的基线差。
\item 若固件输出融合角，单独记录恢复时间、过冲和稳态偏差；若只有六轴数据，仅能评价运动后的\textbf{六轴零位恢复}，不能称滤波器动态响应已通过。
\item 重复至少 5 轮，并记录所有运动期间通讯丢帧、轴状态和异常停车。
\end{enumerate}
\textbf{计算：} $\Delta\theta_\mathrm{return}=\bar\theta_\mathrm{after}-\bar\theta_\mathrm{before}$；恢复时间使用预先定义的静态容差带。
\par\textbf{判定：} 恢复时间和回零差\pending；没有融合角输出则该部分记为“不适用”。

\testheading{CT-17}{时间对齐与动态响应诊断}{P1／有条件；非绝对延迟认证}
\textbf{目的：} 评价 IMU 六轴数据/倾角输出与转台反馈的相对时序问题，为下一阶段动态频响和融合算法工作准备。
\par\textbf{工况：} 优先选\emph{小幅低速}摆动，例如 $A=5^\circ,f=0.1\,\mathrm{Hz}$（最大速率约 $3.14^\circ/\mathrm s$）；如转台摆动模式与路径未实机核实，则不自动执行。
\par\textbf{步骤：}
\begin{enumerate}
\item 记录 IMU 原始逐样本计数器/设备时戳、主机接收时刻、转台状态发送时刻/主机时戳；核对是否有稳定硬件同一时基。没有时必须标记时钟域不同。
\item 审查正弦幅度-频率对应的最大角速度 $\omega_{\max}=2\pi fA$（$A$ 使用角度时得 $^\circ$/s），并保证小于软件、机械和电缆约束。
\item 在确认安全后完成 3--5 个完整周期；记录转台实际角位置而不是只记录目标波形。
\item 将两流时序转换到单一候选时基，检查丢包、插值、采样间隔、输出时延；相关匹配只作为\textbf{相对时间偏移估计}，因为转台编码器读取、传输与滤波还含未知延迟。
\item 输出补偿前后的波形误差和相位图（如果有可对比的倾角输出）；同时记录数据时间戳不确定度，不能自动用交叉相关对齐后宣称端到端延迟合格。
\end{enumerate}
\textbf{输出：} 相对滞后 $\hat\tau$、波形图、丢帧率、采样周期统计、时间基准说明。
\par\textbf{判定：} 当前\texttt{DIAGNOSTIC}；绝对相位/延迟指标的正式验收属于阶段 S3，并要求独立时钟或已评估不确定度的同步参考。

\testheading{CT-18}{独立测试报告、结果等级及数据封存}{P0／必做归档}
\textbf{目的：} 保证所有 CT 实际条件与结论可追踪、可复算，避免“报告只保留通过项”。
\par\textbf{步骤：}
\begin{enumerate}
\item 审核 CT-00--CT-17（含 CT-12A）的完整目录、测试起止时刻、工况编号、样本数、失败重试与剔除比例，保留所有\texttt{FAIL}/\texttt{BLOCKED} 的原始记录。
\item 检查是否存在共用校准数据重新参与验收、不同 PCB 共用参数文件、测试中改动滤波器/量程/ODR、超过转台工作环境、无速度计量却签署速率精度等问题。
\item 汇总每个测试的参考值、均值、标准差、最大绝对误差、产品限值、测量不确定度与结果等级；限值缺失的项目为\texttt{待审批}，不能“默认通过”。
\item 生成 S1 归档清单：原始六轴、转台反馈、实验事件、冻参文件、工艺配置、脚本版本、报告和图表，并对文件计算 SHA-256。
\item 标明本次能得出的结论范围：“常温条件下 PCB 六轴内参后验证”，以及未完成的速率绝对精度、全温、纯线性加速度、振动、整机安装和动态融合测试。
\end{enumerate}
\textbf{交付：} 《S1 阶段测试报告》+ 表格记录 + 原始数据归档包 + 问题单和整改闭环表。
\par\textbf{判定：} 证据齐全、结论与参考能力一致、正式限值已审批才可签发；其他情况只形成\pilot 的研发报告。
